\chapter{独一无二的个体}\label{ch:g9:unique-individuals}

七十亿人，没有两个相同的——不在同一个家庭，
不在同一个操场，不在整个人类历史中，同卵双胞胎
也只算半个例外。前三章已经组装起这一事实的全部机制。
本章是简短的总结章，把整套论证从头到尾走一遍，
并归纳其推论：关于\omterm{def:g5:biodiversity:species}{物种}的多样性，
关于如何读懂“\omterm{def:g9:genes-and-dna:gene}{基因}”类新闻标题，
以及“独一无二”究竟意味着什么、不意味着什么。

\section{唯一性定理}

\begin{proposition}[为何没有两人完全相同]\label{prop:g9:unique-individuals:theorem}
综合三种机制：
\begin{enumerate}
  \item \emph{减数分裂的洗牌}：每位父母的
        \omterm{def:g8:sexual-reproduction:gametes}{配子}在每对同源\omterm{def:g9:chromosomes:chromosomes}{染色体}中各携一条，
        各自独立发牌——每位父母可产生 $2^{23}$ 种半副牌
        （见\cref{rem:g9:chromosomes:whyhalve}），实际数目远不止于此，
        因为成对的同源\omterm{def:g9:chromosomes:chromosomes}{染色体}在减数分裂中还会\emph{互换
        片段}，
        把切分深入到整条\omterm{def:g9:chromosomes:chromosomes}{染色体}之下；
  \item \emph{\omterm{def:g5:flower-to-fruit:steps}{受精}的抽签}：数以百万计的\omterm{def:g8:reproductive-systems:male}{精子}中，
        哪一个（见\cref{ex:g8:reproductive-systems:sperm}），
        与哪个月的卵子相遇，
        使两者的可能性相乘；
  \item \emph{\omterm{prop:g9:genes-and-dna:explains}{突变}的细流}：每一代都会新增
        一些全新的“拼写”，
        （见\cref{prop:g9:genes-and-dna:explains}）。
\end{enumerate}
组合空间远超历史上存在过的全部人类数量：
每一次\omterm{def:g5:flower-to-fruit:steps}{受精}都会发出一副此前从未出现过的牌，
此后也不会再有相同的一副。
\end{proposition}
\admitted

\begin{example}[双胞胎与例外的界限]\label{ex:g9:unique-individuals:twins}
同卵双胞胎发到的是同一副牌
（见\cref{exo:g8:from-egg-to-birth:10}）——即便如此，他们之间也会
出现分化：分开生活会形成不同的后天\omterm{def:g9:heredity-and-traits:traits}{性状}
（见\cref{prop:g9:heredity-and-traits:sorting}），不同的
\omterm{def:g8:nervous-communication:synapse}{突触}经历
（见\cref{ex:g8:nervous-communication:juggler}），甚至
各自体\omterm{def:g5:first-look-at-cells:cell}{细胞}中还有少量不同的晚期\omterm{prop:g9:genes-and-dna:explains}{突变}。唯一的
独一无二性上的裂缝会自行弥合：发牌之后，
生活会使二者逐渐分化。
\end{example}

\begin{remark}[独一无二，却又大体相同]\label{rem:g9:unique-individuals:alike}
两个事实都要记住：任意两人的 \omterm{def:g9:genes-and-dna:dna}{DNA} 序列
一致率超过 $99$\%——这正是
\cref{prop:g5:first-look-at-cells:principle}所述同一\omterm{def:g5:biodiversity:species}{物种}的
共同结构——而分散在三十亿个碱基中的差异，
仍足以保证严格意义上的独一无二。“全部不同”与
“全部是亲属”，不过是同一事实在两种
放大倍数下的样子
（见\cref{met:g9:genes-and-dna:levels}）；以往拿前一事实去对抗后一事实的做法，
首先是糟糕的生物学，
然后才是糟糕的伦理。
\end{remark}

\section{独一无二性的应用}

\begin{example}[DNA 鉴定]\label{ex:g9:unique-individuals:forensics}
独一无二性是可以检验的，法庭也正是这样做的：只要某人足够多的
可变位点（从一小块\omterm{def:g5:first-look-at-cells:cell}{细胞}痕迹中读出），
便能在整个人类中锁定唯一的个体（同卵双胞胎除外）。
同样的比对若在家族层面进行，
便可判定\omterm{prop:g6:classification-kinship:kinship}{亲缘关系}——孩子的每一个可变位点
都必须能从双亲各自的牌中找到来源。
“相册”的逻辑
（见\cref{ch:g9:heredity-and-traits}）如今已成为一种
工具。
\end{example}

\begin{example}[输血与移植]\label{ex:g9:unique-individuals:medicine}
医学每天都在与独一无二性打交道。血型
（见\cref{ex:g9:genes-and-dna:bloodgroups}）是较为简单的情形：
只分四大类，可依据登记信息进行配型。
器官移植则面临较难的情形：免疫系统
（在\cref{ch:g9:immune-defenses}中讲解了其机制）
能够识别更为精细的“自我标志”，其独特程度之高，
以至于供体首先要在家族内搜寻——
亲缘越近，所发之牌越相近。
\end{example}

\begin{proposition}[多样性是物种的储备]\label{prop:g9:unique-individuals:reserve}
\Cref{rem:g8:sexual-reproduction:reserve}曾许诺过洗牌的意义，
如今这些机制兑现了这一许诺。由独一无二的个体组成的\omterm{def:g5:biodiversity:species}{物种}，
其中散布着各种应对情形的\omterm{def:g9:genes-and-dna:gene}{等位基因}，
有的个体携带针对这种
疾病的抗性，有的个体携带对那种胁迫的耐受力，
（即\cref{ex:g8:sexual-reproduction:bets}中的果园，
放到\omterm{def:g5:biodiversity:species}{物种}尺度上即是如此）。\omterm{def:g6:exploring-our-environment:components}{环境}\omterm{def:g6:exploring-our-environment:conditions}{条件}一旦改变，
\cref{prop:g8:reproduction-and-environment:limits}所述的制约
便会不均匀地落在这些差异之上——而这种不均匀
一代又一代所产生的作用，正是
接下来两章要展开讨论的主题：\omterm{def:g5:biodiversity:species}{物种}的
演化。
\end{proposition}
\admitted

\begin{method}[审核“基因”相关的说法]\label{met:g9:unique-individuals:claims}
对于任何关于\omterm{def:g9:genes-and-dna:gene}{基因}与人的标题：
\begin{enumerate}
  \item 层次检验（见\cref{met:g9:genes-and-dna:levels}）：
        这一说法针对的是一个\omterm{def:g9:genes-and-dna:gene}{基因}、一副牌，
        还是家族中的\omterm{def:g9:heredity-and-traits:heredity}{遗传}模式？
  \item 交织检验
        （见\cref{prop:g9:heredity-and-traits:sorting}）：它是否
        尊重“\omterm{def:g9:heredity-and-traits:heredity}{遗传}范围与生活交织”的图景，
        还是假装一个\omterm{def:g9:genes-and-dna:gene}{基因}就决定一生的命运？
  \item 多样性检验：它是否把\omterm{def:g8:reproduction-and-environment:population}{种群}的分布
        视为它本身的储备，还是视为对
        一副并不存在的“正常”牌的偏离？
  \item 双胞胎检验：这一说法能否经得起同卵双胞胎的检验，
        他们共享同一副牌，却依然各不相同？
\end{enumerate}
\end{method}

\section{练习}

\begin{exercise}[$\star$]\label{exo:g9:unique-individuals:1}
列出唯一性定理的三种机制。
\end{exercise}

\begin{exercise}[$\star$]\label{exo:g9:unique-individuals:2}
减数分裂时的片段互换，给洗牌
增加了什么？
\end{exercise}

\begin{exercise}[$\star$]\label{exo:g9:unique-individuals:3}
用一句话同时说出\cref{rem:g9:unique-individuals:alike}
的两个要点——百分比与唯一性——
。
\end{exercise}

\begin{exercise}[$\star$]\label{exo:g9:unique-individuals:4}
同卵双胞胎究竟如何仍能被区分？
\end{exercise}

\begin{exercise}[$\star$]\label{exo:g9:unique-individuals:5}
\omterm{def:g9:genes-and-dna:dna}{DNA} 比对能为法庭
回答哪两个问题？
\end{exercise}

\begin{exercise}[$\star$]\label{exo:g9:unique-individuals:6}
为什么移植施主首先要在家族成员中寻找？
\end{exercise}

\begin{exercise}[$\star\star$]\label{exo:g9:unique-individuals:7}
推广这一论证：把双亲各自的
半副牌空间与\omterm{def:g5:flower-to-fruit:steps}{受精}的“彩票”合并，
得出“此前从未被发出”的结论。
\end{exercise}

\begin{exercise}[$\star\star$]\label{exo:g9:unique-individuals:8}
把\cref{prop:g9:unique-individuals:reserve}应用于
\cref{exo:g8:sexual-reproduction:10}中的葡萄园：
克隆体缺少什么，而由\omterm{def:g2:from-seed-to-plant:seed}{种子}长成的\omterm{def:g8:reproduction-and-environment:population}{种群}却拥有？
\end{exercise}

\begin{exercise}[$\star\star$]\label{exo:g9:unique-individuals:9}
用\cref{met:g9:unique-individuals:claims}检验：
“发现了与学业成功有关的\omterm{def:g9:genes-and-dna:gene}{基因}”。
\end{exercise}

\begin{exercise}[$\star\star$]\label{exo:g9:unique-individuals:10}
为什么“正常的人类一副牌”是
一个没有所指的说法？请结合储备与定理回答。
\end{exercise}

\begin{exercise}[$\star\star$]\label{exo:g9:unique-individuals:11}
一项亲子鉴定排除了一位被怀疑的父亲，
此人曾被\cref{exo:g9:heredity-and-traits:11}中那位邻居错误指控。
请逐个位点说明该检验检查了什么。
\end{exercise}

\begin{exercise}[$\star\star\star$]\label{exo:g9:unique-individuals:12}
“个个各不相同，人人皆有亲缘，两者
皆由同一机制所致。”撰写结语段落：洗牌、
\omterm{prop:g9:genes-and-dna:explains}{突变}与共享密码，各就其位。
\end{exercise}

\section{问题：悬案研讨会}

\begin{problem}[{周末问题——一根头发、三个
问题、三十年}]\label{pb:g9:unique-individuals:1}
一场（虚构的）悬案研讨会：一桩三十年的
未破入室盗窃案，一根保存下来的头发及其毛根
\omterm{def:g5:first-look-at-cells:cell}{细胞}，以及现代比对方法。研讨会引导
学生厘清 \omterm{def:g9:genes-and-dna:dna}{DNA} 鉴定能说明什么、
不能说明什么。

\medskip\noindent\textbf{第一部分——头发
所蕴含的信息。}
\begin{enumerate}
  \item 为何必须保存头发的\emph{毛根}，
        才能完整读取（\cref{def:g6:cell-unit-of-life:parts}
        知道 \omterm{def:g9:genes-and-dna:dna}{DNA} 位于何处）？
  \item 实验室读取一组变异性很高的
        位点。为何选择变异大的位点？
        在 99\% 共享的位点上能确立什么？
  \item 图谱在每个位点上都与一名嫌疑人吻合。
        请说明这一吻合断言了什么，以及它唯一的
        系统性例外
        （\cref{ex:g9:unique-individuals:twins}）。
  \item 辩方指出，嫌疑人在国外有一位同卵
        双胞胎。法庭现在必须做什么，以及
        为什么？
\end{enumerate}

\medskip\noindent\textbf{第二部分——家族
相关问题。}
\begin{enumerate}[resume]
  \item 较早前，一份不完整的图谱曾指向“
        K 家族的一位近亲”。请解释
        \emph{部分}匹配如何读作\omterm{prop:g6:classification-kinship:kinship}{亲缘关系}
        （\cref{prop:g9:unique-individuals:theorem}
        中的发牌过程，倒着运行）。
  \item 一位陪审员问，图谱是否揭示嫌疑人的
        特征——外貌、性情，
        “犯罪\omterm{def:g9:genes-and-dna:gene}{基因}”。请用
        \cref{met:g9:unique-individuals:claims}
        逐项检验这一问题。
  \item 另一位陪审员问，这一吻合可能是“本城一例”
        而非“全人类一例”。
        什么决定了读取位点的
        数量？
  \item 研讨会主持人提醒：“生物学说明是
        谁的\omterm{def:g5:first-look-at-cells:cell}{细胞}；案件仍须说明它们
        如何到达那里。”请区分这两个论断。
\end{enumerate}

\medskip\noindent\textbf{第三部分——研讨会的
结语。}
\begin{enumerate}[resume]
  \item 用一句话概括这一工具：说明
        两章的机制
        （\cref{ch:g9:chromosomes}，
        \cref{ch:g9:genes-and-dna}）
        如何使图谱唯一、家族关系可读。
  \item 为何同一工具也适用于一根三十年前的头发？
        （所依赖的是 \omterm{def:g9:genes-and-dna:dna}{DNA} 的何种性质——
        除\cref{prop:g9:genes-and-dna:explains}所说的复制之外，
        被依赖的是稳定性。）
  \item 研讨会以伦理问题收尾：列出本章所示的
        唯一性服务于人的两种用途，以及
        历史上的误用，
        \cref{rem:g9:unique-individuals:alike}所警示
        的情形。
  \item 以为法庭重新表述的定理作结：
        什么从未发生过两次，以及
        原因何在。
\end{enumerate}
\end{problem}
